% Literature synthesis prepared 17 September 2026. No experimental results are asserted here.
\chapter{Related Work}
\label{chp:related_work}

Research on process mining and simulation addresses two connected problems: reconstructing how a process operates from recorded events, and using an executable representation to assess how its performance might change. Their connection is particularly relevant in logistics, where routing decisions, physical movement and competition for shared resources jointly determine performance. Nevertheless, a model that reproduces recorded activity sequences does not necessarily reproduce the mechanisms that generate waiting times or the consequences of changing resource capacity. The distinction between these objectives provides the organising principle for this chapter.

Building on the concepts introduced in Chapter~\ref{chp:background}, the chapter compares research on process discovery, data-aware simulation, and simulation in logistics. It then examines object-centric and resource-centric alternatives and the evaluation challenges that remain across these approaches. The review is a focused narrative synthesis of foundational methods, relevant application studies and recent methodological work, including publications available in September 2026. It does not claim the exhaustive coverage or screening protocol of a systematic review. The emphasis is on the assumptions, evidence and transferability of each research direction to truck-visit simulation at Container Terminal Burchardkai (CTB).

\section{Process Discovery as a Basis for Simulation}
\label{sec:rw_discovery}

\subsection{Discovery Algorithms and Their Representational Choices}

Automated process discovery constructs a process model from an event log, but the choice of discovery technique determines which structures can be represented and how infrequent behavior is handled. Augusto et al. compare discovery approaches along several quality dimensions and demonstrate why a single measure cannot adequately characterise their output~\cite{RWAugusto2019Review}. A highly permissive model may reproduce most observed traces while allowing substantial additional behavior; a restrictive model may capture frequent variants while excluding legitimate exceptions. For simulation, this trade-off affects the behavior that the execution engine can subsequently generate.

Inductive Miner addresses discovery through recursive decomposition into structured process fragments. Its process-tree representation provides a basis for constructing sound executable models, while extensions for infrequent behavior introduce filtering choices~\cite{Leemans2014InductiveMiner}. This makes the approach a useful baseline, but its structural guarantees do not establish that the discovered ordering or concurrency matches the physical process. Heuristics-based discovery instead uses frequency and dependency information to identify relations between activities. Region-based approaches, including integer-linear-programming methods, construct places subject to constraints derived from observed behavior. These families offer different assumptions and failure modes rather than interchangeable implementations of the same principle~\cite{RWAugusto2019Review,RWVanZelst2018ILP}.

Split Miner provides a further comparator by combining filtering and concurrency identification to construct relatively simple BPMN models while balancing fitness and precision~\cite{RWSplitMiner2019}. Its results in benchmark settings motivate consideration, but do not establish it as the best algorithm for logistics. A simulation pipeline that requires Petri nets must also account for the conversion from the discovered representation to the executable one. Guarantees concerning a particular algorithm, model class or conversion cannot be transferred automatically to a different software variant. This is especially relevant for ILP-based methods: relaxed soundness and full workflow-net soundness are different properties, and an implementation labelled ``ILP Miner'' need not reproduce every property of a published variant~\cite{RWVanZelst2018ILP}.

The methodological question is therefore which configurations produce an appropriate model for a defined log and intended use. Rehse et al. argue that discovery experiments require explicit assumptions and reproducible experimental designs; results obtained with different logs, settings and metrics cannot simply be combined into a universal ranking~\cite{RWRehse2024Experimentation}. For CTB, a controlled comparison must hold the input abstraction and evaluation cohort constant, document filtering and computational failures, and assess important route variants separately. Such a comparison contributes evidence for model selection without requiring a new discovery algorithm.

\subsection{Lifecycle Information, Rules and Explained Deviations}

The semantics of recorded events constrain what discovery can establish. Leemans et al. distinguish concurrency from interleaving: activities may occur in either order without being able to execute simultaneously. Their lifecycle-aware discovery approach uses activity starts and completions to expose this distinction and improve the interpretation of performance information~\cite{RWLeemans2016Lifecycle}. Split Miner 2 similarly uses execution intervals to identify true concurrency and distinguish exclusive from inclusive choices~\cite{RWSplitMiner2Preprint}. The implication for logistics is substantial. A truck may visit two handling areas in different orders across visits, yet cannot necessarily receive two physical services at the same time. Completion-only observations cannot resolve this question, while readiness or assignment timestamps cannot be treated as service starts without additional evidence.

Domain rules provide another source of information beyond event frequencies. Norouzifar et al. introduce rule-guided process discovery, integrating discovered or externally specified constraints into the recursive selection of process structures~\cite{RWNorouzifar2026Rules}. Their evaluation considers both conventional conformance measures and agreement with rules. It also shows that a dominant variant can conceal deficiencies affecting less frequent behavior. However, rule guidance does not eliminate representational restrictions, including difficulties with long-term dependencies. Rules extracted from the same log also have a different evidential status from independently documented operational requirements.

Conformance checking complements discovery by identifying discrepancies between model and log~\cite{aalst2012ConformanceChecking}. Mozafari Mehr et al. extend this perspective with contextual explanations of deviations involving control flow, data operations and authorisation~\cite{RWMehr2023Explainable}. Their multi-layer approach is evaluated in a financial process rather than logistics, but it establishes a useful analytical principle: the meaning of a deviation depends on its context. A missing modeled activity may reflect an incomplete record, an overly restrictive model or a genuine exception. For a logistics case study, reporting aggregate fitness should therefore be accompanied by an examination of affected routes and activities. A simplified contextual analysis should not be presented as an implementation of the paper's complete conformance method.

\section{From Discovered Processes to Data-Aware Simulation}
\label{sec:rw_data_aware}

\subsection{Data-Driven Construction and Stochastic Behavior}

Rozinat et al. establish an early connection between event logs and executable simulation models by extracting multiple process perspectives~\cite{PMRozinat2009}. Their work moves beyond discovering an activity sequence: simulation also requires timing, resource and decision information. Camargo et al. subsequently address the calibration effort through an accuracy-oriented discovery procedure that searches over configuration choices and evaluates similarity between observed and simulated behavior~\cite{Camargo2020}. The more recent SIMOD implementation organises discovery into a configurable pipeline rather than a single control-flow mining operation~\cite{simod2024}.

These approaches illustrate the distinction between \emph{data-driven} and \emph{data-aware} simulation. A model is data-driven when its structure or parameters are inferred from observations. It is data-aware when data values influence execution, for example through routing conditions or attribute-dependent timing. An automatically fitted model using fixed branch probabilities and unconditional duration distributions is data-driven, but does not necessarily express how case attributes change decisions~\cite{lopez2024}.

Stochastic process mining further distinguishes the behavior permitted by a model from its likelihood. Rogge-Solti et al. investigate the discovery of stochastic Petri nets with general delay distributions, while Burke et al. study the estimation of stochastic transition weights~\cite{stochastikPetriNets2013Solti,burke2020stochasticPD}. Both perspectives matter for simulation: allowing a route is insufficient if its generated frequency or duration is unrealistic. They nevertheless address different components of the simulation problem. Reproducing a trace distribution does not, by itself, identify the resource and transport mechanisms that produced it.

\subsection{Data Conditions, Resources and Runtime Prediction}

Data-aware approaches differ in what they model explicitly. Mannhardt et al. extend stochastic process models with data-dependent transition weights and corresponding conformance analysis~\cite{mannhardt2023stochasticProcessesModelling}. Related work investigates how process-state information affects activity-firing probabilities~\cite{Vinci2023Inf.Activ.Prob}. These approaches can distinguish situations that share the same control-flow position but differ in their attributes or execution history. This is relevant when the next truck operation depends on the remaining tasks rather than only on the currently enabled transitions.

L\'opez-Pintado et al. address an additional requirement: simulation must specify how the data used by decisions are generated and updated. Their data-aware discovery approach distinguishes global, case-level and event-level attributes and supports deterministic and stochastic update rules together with branching conditions~\cite{lopez2024}. This is more expressive than conditioning decisions on a set of externally fixed case attributes. Its benefit depends on whether recorded attribute changes can be associated with the correct execution steps. An attribute observed only after a task finishes cannot automatically serve as an input to the decision that started it.

The resource perspective introduces a separate source of heterogeneity. L\'opez-Pintado and Dumas show that treating resources as interchangeable can obscure differences in availability and performance. Their differentiated-resource models more closely reproduce cycle-time and work-rhythm distributions in the evaluated logs~\cite{BPSdifferentiatedResources2022}. Prosimos implements this line of work through a simulation engine with explicit resource allocation and calendars~\cite{Prosimos2023}. These findings motivate differentiated resource models where the required evidence exists, but a yard-area identifier is not necessarily an individual crane identifier. The level of differentiation must match the granularity of the observations.

Hybrid methods use learned predictors to represent dependencies that are difficult to capture through fixed parameters. Camargo et al. compare data-driven simulation with deep-learning-based generation, exposing differences in accuracy and the ability to represent process behavior~\cite{Camargo2021}. Meneghello et al. integrate prediction with simulation at runtime in RIMS~\cite{runtimeIntegration2024}. This allows predictions to use features calculated from the evolving simulation state, such as current workload and queue information. It avoids the inconsistency of adding context-dependent times only after an entire event sequence has already been generated. However, an executable control-flow model does not make every learned decision transparent, and the behavior of predictors under substantially changed operating conditions remains an evaluation question.

\subsection{Probabilistic White-Box Simulation}

ProSiT adopts probabilistic decision trees to characterise several runtime perspectives while retaining inspectable rules~\cite{Vinci2026prosit}. The approach accepts a Petri-net control-flow model and learns conditional transition weights and temporal behavior. Tree leaves represent probabilities or fitted distributions, enabling the modeler to inspect how attributes and process history influence execution. This creates a direct connection between representation and scenario construction: changing a condition or leaf parameter corresponds to an explicit change in the runtime model.

The evaluation of ProSiT holds the control-flow backbone constant when comparing simulation approaches and assesses multiple event-log distances, decision-level behavior and variability~\cite{Vinci2026prosit}. This experimental choice is important because changing the net and the runtime predictor simultaneously would obscure which component explains an improvement. The paper nevertheless does not establish that ProSiT is the most suitable framework for every logistics setting. Its waiting-time learning depends on transition enablement and aligned event data. Consequently, the interpretation of the control-flow model affects the temporal targets learned from the log.

White-box representation provides inspectability, but not automatic operational correctness. A shallow tree may be understandable while depending on a poorly defined timestamp or a variable unavailable at decision time. Moreover, waiting represented through a learned distribution must be distinguished from waiting generated by explicit resource competition. Combining these mechanisms requires a precise account of which delay each component represents. This is a central consideration when enriching a data-aware simulator with fixed transport times or capacity constraints.

\section{Process Mining and Simulation in Logistics}
\label{sec:rw_logistics}

\subsection{Logistics Applications and the Evidence Base}

Alnahas reviews applications of process mining to logistics in manufacturing organisations, identifying fifteen studies from a search covering 2004--2022~\cite{RWAlnahas2023Logistics}. The selected studies illustrate substantial use of discovery, with a smaller emphasis on combinations of discovery, conformance and enhancement. The review establishes that process mining has been applied to logistics, but its scope and search period do not support claims about all current logistics applications. In particular, the existence of a discovered process model should not be interpreted as evidence of an executable, validated simulation model.

Becker and Intoyoad investigate context-aware process mining in logistics, addressing the heterogeneous systems and contextual information involved in reconstructing logistics processes~\cite{RWBecker2017Context}. Their work connects data preparation to the interpretation of operational behavior. For a terminal, combining gate records, handling orders and equipment events requires more than matching column names: timestamps and identifiers may describe different operational stages. Event-log construction therefore introduces modeling choices before any discovery algorithm is applied.

The 2025 study by Grobis and Ihlenfeldt outlines how process mining can support data preparation and simulation design in production and logistics~\cite{Grobis2025ProductionLogistics}. It emphasises semantic consistency, resource relationships and the integration of information at different levels of detail. Its contribution is primarily a proposed approach and requirements analysis; practical validation is identified as subsequent work. Their 2026 production-planning case then demonstrates a concrete transformation from event data into simulation-relevant structures and parameters~\cite{Grobis2026SimulationDevelopment}. These papers should be read as successive contributions with different forms of evidence, rather than as two independent demonstrations of general effectiveness.

Healthcare logistics offers adjacent evidence but also requires careful delimitation. Murazzano et al. describe a systematic review \emph{protocol} covering applications such as transport, inventory and scheduling~\cite{RWMurazzano2026Protocol}. The publication motivates an operational-logistics perspective beyond clinical pathways, but does not report the results of a completed review. By contrast, Vinci et al. present an implemented healthcare case study connecting model discovery, simulation assessment and resource-allocation scenarios~\cite{Hospital2026Vinci}. Its methodological sequence is transferable to CTB, whereas its resource assumptions and reported improvements are specific to the hospital setting.

\subsection{Terminal Simulation and Operational Decisions}

Container-terminal simulation has a substantial history independent of process mining. Yun and Choi use an object-oriented simulation approach to represent terminal operations, while Hartmann studies the generation of scenarios for terminal simulation and optimisation~\cite{YUN1999DESSimulation,Hartmann2005cenariosforCTSimulation}. Bruzzone et al. address the simulation, analysis and optimisation of terminal processes~\cite{Bruzzone2012}. These works establish the importance of executable resource and material-flow representations. They do not imply that these representations can be recovered automatically from transactional event logs. Object-oriented software design also differs from object-centric process mining, which preserves relationships between multiple objects in event data.

More recent terminal research connects simulation to explicit operational objectives. Kastner et al. compare simulation-based optimisation procedures for equipment quantities and operational policies~\cite{Kastner2021}. Caballini et al. combine clustering and optimisation to match container transactions and assign truck appointments~\cite{CTprocessMiningTAS2020}. The latter is a data-mining and optimisation approach, not a Petri-net process-discovery study. Bett et al. investigate truck appointments and dual transactions through discrete-event simulation~\cite{Bett2024}. Together, these studies show that the relevant output is a justified operational comparison, rather than a process diagram alone.

The connection to data-aware process simulation remains an integration problem. Terminal studies provide detailed decision objectives and physical constraints, whereas general simulation-discovery methods provide ways to infer behavior from logs. Combining them requires establishing which mechanisms are observable and which must be supplied from other evidence. A model may reproduce average truck turnaround time while misrepresenting the balance between movement, service and waiting. Such compensation becomes problematic when a scenario changes one of those components.

\subsection{Transport, Waiting and Partially Observed Mechanisms}

Grobis and Ihlenfeldt explicitly decompose an interval between production activities into transport and waiting~\cite{Grobis2026SimulationDevelopment}. In their case, transport is represented as a deterministic delay, while waiting emerges from resource and buffer interactions. This provides a concrete precedent for distinguishing physical movement from congestion. It does not justify transferring a constant duration to truck operations without route-specific evidence. The decomposition depends on the meaning of the recorded boundary events and on observations or assumptions that distinguish transport from other delays.

Related business-process research identifies \emph{extraneous delays}: delays not explained by resource contention or unavailability. Chapela-Campa and Dumas propose discovering such delays and representing them explicitly in a simulation model~\cite{RWChapela2024Delays}. This broadens the treatment of elapsed time, but an inferred residual delay is not automatically a measured transport duration. For CTB, the correspondence between residual time and a physical mechanism must be argued separately.

The main methodological risk is double counting. If a fitted duration already contains movement or waiting, adding a transport delay or resource queue can count the same elapsed time again. Conversely, replacing all observed waiting with a fixed delay removes the response to workload that an intervention study may need. A defensible enrichment therefore requires a consistent decomposition of time and a sensitivity analysis when its components cannot be uniquely identified from the available observations.

\section{Object-Centric and Resource-Centric Alternatives}
\label{sec:rw_alternatives}

Traditional event logs organise events around one case identifier. Object-centric process mining instead allows events to refer to multiple related objects, addressing the loss of relationships caused by flattening operational records into independent traces~\cite{VanDerAalst2023OCPM}. In logistics, a handling event may concern a truck visit, a container, a task and a resource. Selecting only one perspective can duplicate events or conceal dependencies. However, adding attributes to a truck-visit log does not by itself make the model object-centric.

Knopp et al. develop object-centric simulation discovery, addressing both the generation of related objects and their synchronisation and routing through a model~\cite{Knopp2023ObjectCentricSimulation}. This directly extends the single-case simulation perspective. The additional expressiveness also introduces requirements: object relationships must be present or reconstructable, and simulation must preserve the relevant identities during execution. The existence of an object-centric process model is therefore not sufficient to establish correct resource allocation or spatial movement.

Resource-centric simulation offers a different alternative. Kirchdorfer et al.'s AgentSimulator discovers interacting agents representing resources and their behavior~\cite{RWKirchdorfer2025Agents}. This addresses settings where individual decisions and handovers shape the process rather than merely execute a centrally prescribed route. Resource-centricity and object-centricity answer different questions: one emphasises who acts and interacts; the other preserves which objects participate in events. Neither is synonymous with data awareness, although both may incorporate data-dependent behavior.

For CTB, these alternatives identify the limits of the chosen abstraction. If the research question concerns truck turnaround and visit-level routing, a visit-centric model with explicit shared-resource logic may be adequate. If the question depends on synchronising several containers, tasks or vehicles, a single-case representation may become insufficient. The choice should follow the dependencies needed for the operational question and the evidence available to represent them, rather than the recency of a modeling paradigm.

\section{Validation and Current Research Challenges}
\label{sec:rw_validation}

\subsection{Model Quality, Generated Behavior and Intended Use}

Sulis et al.'s review identifies validation as an unresolved issue across the intersection of simulation and process mining~\cite{Sulis2026SimulationReview}. Of the 105 reviewed studies, 36 do not report validation, while qualitative assessments often provide insufficient detail about their procedure. The review also highlights limited support for queues and resource behavior. These findings support distinguishing three claims: that a model is executable, that it reproduces relevant observations, and that it is suitable for a specified decision. The broader simulation literature makes the intended use central to verification and validation~\cite{Sargent2013Validation,BPSAalst2015}.

Conformance metrics primarily compare recorded behavior with the behavior permitted by a process model. Stochastic conformance additionally considers behavioral frequencies, as illustrated by Earth Movers' Stochastic Conformance Checking~\cite{leeman2013EarthMovers}. Simulation evaluation must go further because the execution engine generates timestamps, routes and resource interactions. Chapela-Campa et al. propose a multi-perspective framework for comparing real and simulated event logs, including control-flow, temporal and congestion-related characteristics~\cite{Chapela2025SimQuality}. The hospital study applies such comparisons alongside process-model assessment, repeated simulations and subgroup analysis~\cite{Hospital2026Vinci}.

These metrics are complementary rather than interchangeable. Similar cycle-time distributions can coexist with incorrect route frequencies, and aggregate agreement can conceal poor behavior for a minority operation type. The unit of analysis, normalisation, observation window and subgroup sizes therefore affect the interpretation of results. Simulation replications quantify variability conditional on a model; they do not alone quantify uncertainty in reconstructed timestamps, fitted parameters or structural assumptions.

\subsection{Temporal Transfer, Utility and Operational Forecasting}

Recent work challenges the assumption that similarity to a later event log is always the appropriate target. In their utility-based evaluation proposal, \"Ozdemir et al. distinguish reproducing an observed process from forecasting a future period that may contain drift~\cite{Ozdemir2025UtilityBPS}. They evaluate whether simulated data can support downstream predictive tasks comparably to real data. This supplies a different quality signal and draws attention to temporal patterns that distributional summaries may conceal. However, success in predictive-model training is not itself evidence that a terminal intervention is represented correctly.

A chronological holdout remains useful for evaluating transfer to a later operating period, but the claim must be named accordingly. If arrival patterns or the mix of operations change, an observed discrepancy may reflect changing conditions as well as model error. Retrospective reconstruction, later-period transfer and operational scenario analysis should consequently be reported separately. This distinction is particularly relevant to logistics, where changes in demand and resource deployment can alter congestion without changing the set of activity labels.

Avramenko et al. extend short-term simulation through reconstruction of the current process state and uncertainty quantification~\cite{RWAvramenko2026ShortTerm}. Their work shows why starting an operational simulation requires attention to already active cases and workload rather than an empty system alone. For a terminal study, this distinguishes tactical comparisons of recurring operating conditions from forecasts of the next few hours. The latter introduce additional requirements for state information and forecast uncertainty; historical model fit does not remove them.

\subsection{Explainability, Data Quality and Intervention Assumptions}

Explainability has both representational and empirical aspects. ProSiT makes learned rules inspectable, while explainable conformance checking connects deviations to their context~\cite{Vinci2026prosit,RWMehr2023Explainable}. Neither substitutes for checking whether the underlying observations describe the intended activities. Similarly, privacy-preserving event-log transformation, as studied in PRETSA, introduces a trade-off between protecting information and retaining useful process behavior~\cite{FahrenkrogPetersen2019PRETSA}. Any transformation that changes case composition or timing must be considered when interpreting discovery and simulation results.

The remaining challenge is the response to intervention. A learned association between workload and waiting can reproduce historical patterns without identifying how waiting would change after adding a crane or modifying an appointment policy. Scenario analysis therefore requires explicit assumptions about which mechanisms remain stable. The purpose of domain review is to assess those assumptions and identify omitted constraints, not to replace quantitative evidence. Its provenance and procedure must be documented, and author-led assessment must remain distinguishable from independent stakeholder validation~\cite{Sulis2026SimulationReview,Sargent2013Validation}.

\section{Synthesis and Positioning of This Thesis}
\label{sec:rw_positioning}

Table~\ref{tab:rw_comparison} summarises the complementary contributions of the reviewed research. The comparison concerns the cited approaches and their reported scope; it is not an exhaustive feature comparison of current software versions.

\begin{table}[htbp]
\centering
\begin{tabularx}{\textwidth}{@{}p{0.23\textwidth}YY@{}}
\hline
Research direction & Main contribution & Question remaining for CTB \\
\hline
Discovery benchmarks and rules~\cite{RWAugusto2019Review,RWNorouzifar2026Rules} & Compare permitted behavior and incorporate additional constraints. & Does the chosen net preserve relevant routes and defensible execution semantics? \\
Data-aware simulation~\cite{lopez2024,Vinci2026prosit} & Represent attribute updates or conditional runtime behavior. & Are features available at decision time, and do learned delays have the intended meaning? \\
Differentiated resources and runtime integration~\cite{BPSdifferentiatedResources2022,runtimeIntegration2024} & Represent heterogeneity and dependencies on evolving workload. & Can the recorded resource information support the required capacity and queue mechanisms? \\
Logistics simulation development~\cite{Grobis2026SimulationDevelopment,Kastner2021} & Connect physical mechanisms and operational objectives to simulation. & Which timing and routing constraints can be recovered from CTB records? \\
Object- and resource-centric simulation~\cite{Knopp2023ObjectCentricSimulation,RWKirchdorfer2025Agents} & Preserve interacting objects or resource-specific behavior. & Does the decision question require a richer abstraction than a truck visit? \\
Simulation evaluation~\cite{Chapela2025SimQuality,Ozdemir2025UtilityBPS} & Assess generated behavior from several perspectives and intended uses. & Which evidence supports reconstruction, temporal transfer and intervention analysis respectively? \\
\hline
\end{tabularx}
\caption{Research contributions and their implications for the CTB case study.}
\label{tab:rw_comparison}
\end{table}

The reviewed literature does not support claiming that process mining in logistics or data-aware simulation is new. It instead motivates an empirical integration question: how discovery choices and explicitly justified logistics constraints affect the adequacy of a data-aware simulation under the information limitations of a terminal event log. Addressing this question requires linking model-selection evidence to the behavior of the resulting simulator, while making external assumptions visible.

The thesis is positioned as a reproducible case study of that connection. Its discovery comparison should establish the trade-offs between behavioral coverage, precision, structural properties and domain plausibility. Its simulation evaluation should then distinguish the effect of the selected control-flow model from the effect of timing and resource enrichment. Operational scenarios can build on this evidence only within the validated scope. The contribution lies in the documented comparison, the justified integration of process and logistics information, and the assessment of where that integration succeeds or remains underdetermined. It does not depend on proposing a new mining algorithm or presuming that a particular framework will outperform the alternatives.
